\documentclass[11pt,a4paper]{article}
\usepackage[T1]{fontenc}
\usepackage[utf8]{inputenc}
\usepackage[brazilian]{babel}
\usepackage{lmodern}
\usepackage[margin=2.2cm]{geometry}
\usepackage{enumitem}
\usepackage{booktabs}
\usepackage{xcolor}
\usepackage{tcolorbox}
\usepackage{fancyhdr}
\usepackage{amssymb}

\pagestyle{fancy}
\fancyhf{}
\lhead{\small Roteiro de entrevista com fornecedor}
\rhead{\small TCC, César Hoffmann, UNISINOS}
\cfoot{\small\thepage}

\setlength{\parindent}{0pt}
\setlength{\parskip}{0.5em}

\newcommand{\aux}[1]{\par\vspace{0.15em}{\small\textcolor{gray}{\textit{Aprofundar: #1}}}}
\newcommand{\bloco}[2]{\section*{#1 \hfill{\normalsize\textcolor{gray}{#2}}}}
\newcommand{\caixa}{$\square$\ }

\setlist[enumerate]{label=\textbf{\arabic*.}, leftmargin=*, itemsep=0.8em}

\begin{document}

\begin{center}
  {\Large\bfseries Roteiro de entrevista com fornecedor de cardápio digital}\\[0.4em]
  Cardápios digitais e autoatendimento em bares e restaurantes de Porto Alegre\\[0.2em]
  {\small César Hoffmann \textbullet{} orientação: Prof.\ Dr.\ Guilherme Lacerda \textbullet{} Ciência da Computação, UNISINOS}
\end{center}

\begin{tcolorbox}[colback=gray!6,colframe=gray!50,title=Antes de começar,fonttitle=\bfseries]
\caixa Gravação testada (no Meet, conferir em que conta do Google a gravação e a transcrição ficam salvas).\\
\caixa TCLE enviado e aceito antes de começar a gravar.\\
\caixa Termo de Confidencialidade oferecido: o fornecedor deve falar de clientes, preços e arquitetura.\\
\caixa Combinar se o nome da empresa e do produto podem aparecer no trabalho ou se prefere pseudônimo.\\
\caixa Ao começar a gravar: dizer data, empresa e papel do entrevistado e pedir o ``ok'' verbal de consentimento.
\end{tcolorbox}

\textbf{Duração-alvo:} 30 minutos. \textbf{Objetivo:} ouvir quem projeta e vende a solução para entender as decisões de arquitetura, os requisitos que os estabelecimentos pedem e o que o fornecedor observa no mercado de Porto Alegre. O fornecedor não é um caso do estudo; ele é um informante complementar aos estabelecimentos e aos frequentadores.

\textbf{Como conduzir.} Os Blocos 2 e 4 são os mais importantes para o TCC. Em 30 minutos não dá para fazer todas as perguntas: nos Blocos 1, 3 e 5, escolha uma ou duas, e deixe os aprofundamentos em cinza para quando sobrar tempo. Peça exemplos concretos (``lembra de um cliente em que\ldots?''). Não cite números nem relatos do questionário ou das outras entrevistas, para não induzir respostas.

%-------------------------------------------------------------
\bloco{Bloco 0, Caracterização}{3 min}
\begin{enumerate}[series=perguntas]
  \item Conte brevemente a história da empresa e do produto: quando surgiu, que problema queria resolver, onde atua.
  \item Qual é o seu papel na empresa e há quanto tempo você está nessa função?
  \item Mais ou menos quantos estabelecimentos usam a solução hoje, e quantos em Porto Alegre ou no RS? Qual é o perfil típico (porte, tipo de casa)?
\end{enumerate}

%-------------------------------------------------------------
\bloco{Bloco 1, Produto e modalidades}{3 min}
\begin{enumerate}[resume=perguntas]
  \item Que modalidades vocês oferecem (QR Code no celular do cliente, tablet, totem) e qual é a mais vendida? Por quê?
  \item O cliente final só consulta o cardápio ou também faz o pedido e paga pelo sistema? Qual é a proporção entre esses usos?
  \item O cardápio é uma página web, um aplicativo ou um arquivo (PDF)? Por que essa escolha?
\end{enumerate}

%-------------------------------------------------------------
\bloco{Bloco 2, Arquitetura}{9 min}
\begin{enumerate}[resume=perguntas]
  \item Em linhas gerais, como a solução é construída e hospedada? É uma plataforma única que atende todos os estabelecimentos (multi-inquilino) ou cada cliente tem uma instalação própria?
    \aux{como vocês separam os dados de cada estabelecimento? Como fazem para escalar em horários de pico?}
  \item Com que outros sistemas a solução se integra (PDV/caixa, emissão de nota fiscal, estoque, impressora ou tela da cozinha, meios de pagamento, delivery)? Quais integrações dão mais trabalho?
  \item O que acontece quando a internet do estabelecimento ou do cliente cai? Existe algum funcionamento offline ou plano B?
  \item Como o estabelecimento atualiza preços, itens novos e itens esgotados? Isso aparece para o cliente em tempo real?
  \item Que nível de disponibilidade vocês garantem? Já tiveram uma queda grande? O que aprenderam com ela?
  \item O que cada estabelecimento pode personalizar (visual, fotos, categorias, idiomas) e o que é igual para todos?
\end{enumerate}

%-------------------------------------------------------------
\bloco{Bloco 3, Requisitos e evolução}{5 min}
\begin{enumerate}[resume=perguntas]
  \item Quais são as funcionalidades que os estabelecimentos mais pedem? E quais eles quase não usam?
  \item Vocês recebem retorno dos clientes finais (quem está na mesa)? Quais são as reclamações e os elogios mais comuns?
    \aux{como esse retorno chega até vocês: pelo estabelecimento, por avaliações, por dados de uso?}
  \item Já lançaram alguma funcionalidade que não deu certo? Por que não funcionou?
  \item Que dados de uso vocês oferecem ao estabelecimento (itens mais vistos, horários de pico, tempo de pedido)? Os clientes usam esses dados?
\end{enumerate}

%-------------------------------------------------------------
\bloco{Bloco 4, Dados pessoais e acessibilidade}{5 min}
\begin{enumerate}[resume=perguntas]
  \item O cliente final precisa se cadastrar ou informar algum dado (nome, CPF, telefone) para ver o cardápio ou pedir? Quem decide isso: vocês ou o estabelecimento? Para que esses dados são usados?
  \item Na LGPD, como vocês entendem o papel de cada um: o estabelecimento é o controlador e vocês o operador, ou é diferente? Como isso aparece no contrato?
  \item A solução segue alguma diretriz de acessibilidade (por exemplo, WCAG)? Funciona com leitor de tela, permite aumentar a fonte, tem bom contraste?
    \aux{já testaram com pessoas idosas ou com deficiência visual? Algum cliente pediu isso?}
\end{enumerate}

%-------------------------------------------------------------
\bloco{Bloco 5, Mercado e adoção}{4 min}
\begin{enumerate}[resume=perguntas]
  \item Na sua experiência, que tipo de estabelecimento adota mais facilmente e que tipo resiste? Quais são as objeções mais comuns na venda?
  \item Como a pandemia mudou a procura pela solução? E depois dela?
  \item A lei de Porto Alegre de 2024, que exige cardápio físico junto com o digital, afetou vendas, cancelamentos ou o uso do produto?
  \item Como funciona o modelo de cobrança (mensalidade, taxa por pedido, equipamento)? Quando um cliente cancela, quais costumam ser os motivos?
\end{enumerate}

%-------------------------------------------------------------
\bloco{Bloco 6, Encerramento}{1 min}
\begin{enumerate}[resume=perguntas]
  \item Daqui a cinco anos, como você imagina o atendimento em bares e restaurantes?
  \item Tem algo importante sobre esse assunto que eu não perguntei?
  \item Vocês poderiam compartilhar algum dado agregado e anonimizado (por exemplo, proporção de pedidos pelo QR Code ou tempo médio de pedido) para o trabalho?
\end{enumerate}

\begin{tcolorbox}[colback=gray!6,colframe=gray!50,title=Logo depois da entrevista,fonttitle=\bfseries]
\caixa Agradecer e combinar o envio do resumo dos resultados.\\
\caixa Baixar a gravação e a transcrição do Meet e guardar fora de pastas compartilhadas.\\
\caixa Anotar em 5 minutos: impressões gerais, o que mais surpreendeu, contradições com os estabelecimentos e o que não ficou gravado.\\
\caixa Salvar com o nome \texttt{fornecedor\_empresa\_AAAA-MM-DD}.
\end{tcolorbox}

\newpage
\textbf{Anotações}\par
\vspace{0.5em}
\newcount\linhas \linhas=0
\loop\ifnum\linhas<24 \rule{\linewidth}{0.3pt}\par\vspace{0.7em}\advance\linhas by 1 \repeat
\end{document}
